\documentclass[12pt]{letter}
\usepackage[margin=1in]{geometry}
\usepackage{hyperref}

\signature{Aviral Sharma}
\address{Chandigarh University \\ Chandigarh, India \\ aviral98765@gmail.com}

\begin{document}

\begin{letter}{Guest Editors \\ Special Issue: ``Artificial Intelligence Methods for Inverse Problems'' \\ Communications in Nonlinear Science and Numerical Simulation \\ Prof. Davide La Torre, Prof. Herb Kunze, Prof. Oleg Michailovich, Prof. Roman Smirnov}

\opening{Dear Guest Editors,}

We are pleased to submit our manuscript, \textit{``A Physics-Informed Neural
Network Forward Solver for the Rough-Volatility Calibration Inverse Problem:
Dimension Scaling on the Markovian-Lifted PDE and a Fast Warm-Start Hybrid,''}
for consideration in the Special Issue on Artificial Intelligence Methods for
Inverse Problems in \textit{Communications in Nonlinear Science and Numerical
Simulation}.

Our manuscript addresses an inverse problem that is difficult for a reason
specific to this Special Issue's scope: the forward problem itself is
ill-posed before any inversion is attempted. Rough-volatility models describe
market variance with a fractional, non-Markovian kernel, so there is no
finite-dimensional forward PDE to solve, and consequently no forward operator
to embed in a calibration inverse problem at all. We resolve this using a
Markovian lift -- approximating the fractional kernel by $n$ exponentials to
recover a genuine but $(n{+}1)$-dimensional forward PDE -- and build the
resulting forward operator from a physics-informed neural network (PINN).

We believe the manuscript fits the Special Issue's call directly on several
of its named criteria:

\begin{itemize}
\item \textbf{Convergence.} We quantify how the AI forward operator's own
solve-error scales with the lifted dimension $n$ (the forward problem's
governing nonlinear system), showing sub-linear growth out to $n=32$ where a
naive finite-difference alternative is already infeasible by $n=8$.
\item \textbf{Stability.} We identify four training design choices --
a hard-constrained initial-condition ansatz, an integrated-variance baseline
anchor, self-adaptive loss weighting, and a causal time-to-maturity curriculum
-- that are each individually necessary for the forward operator to converge
stably, and we report the ablation showing what fails without each one.
\item \textbf{Generalisation and identifiability.} We extend the forward
operator to take the five model parameters as additional inputs and use it
inside the calibration inverse problem itself, explicitly reporting where its
accuracy generalises and where it does not (the rough, high-vol-of-vol
parameter corners), and diagnosing a training-box identifiability failure that
limits it as a standalone inverse-problem solver.
\item \textbf{Rigorous validation against a same-fidelity reference.} Rather
than validating the AI method against the true (non-lifted) model directly, we
decompose its error against an independent forward solver at the same lifted
dimension, isolating the network's own error from the forward model's
approximation error -- a validation methodology we believe is broadly relevant
to AI methods for inverse problems built on approximated governing equations,
not only to this application.
\end{itemize}

Practically, we show that the trained AI forward operator, used as a
differentiable warm start for a short refinement against the true forward
model, solves the calibration inverse problem for real Bitcoin option data
(Deribit exchange) at approximately nine times fewer expensive forward-model
evaluations than solving it from scratch, while a standalone AI-only solve is
shown honestly to be insufficiently accurate on its own.

This manuscript is original, has not been published previously, and is not
under consideration elsewhere. All authors have approved the manuscript and
agree with its submission to this Special Issue. We have no conflicts of
interest to declare, and the study received no external funding.

Thank you for considering our work. We look forward to your response.

\closing{Sincerely,}

\end{letter}
\end{document}
